% GENERATED FILE. Edit canonical lessons, then run npm run generate:books.
% canonical-source-hash: fnv1a64:3d4b3781370b1234
% canonical-lessons: LA-C520-spiritus, LA-C520-anhelo, LA-C520-singultus, LA-C520-oscito, LA-C520-sternuo, LA-R520-review-513-520

\chapter{Breathing}
\label{ch:la-breathing}
\hlchaptermodality{520}

\begin{chapteropening}
\textbf{By the end of this chapter:} \emph{I can describe panting, sobbing, yawning and sneezing.}

Complete the last lesson of chapter 520: I can describe panting, sobbing, yawning and sneezing.
\end{chapteropening}

\section[spīritus, spīritūs]{spīritus, spīritūs — breath; spirit}

\label{lesson:LA-C520-spiritus}



\begin{quote}
Before the new one: say the Latin for to smell, then the Latin for to perceive.
\end{quote}

\subsection*{You'll want to know: spīritus, spīritūs}
\textbf{spīritus, spīritūs} — \textquotedblleft{}breath; spirit\textquotedblright{}.

\begin{quote}
Post cursum spīritum vix dūcō.
\end{quote}

\begin{quote}
\emph{After the run I can hardly draw breath.}
\end{quote}

Masculine, fourth declension, from \emph{spīrō}. Breath first, then the spirit. English \emph{spirit}.

\begin{cousinweb}
\textbf{spīritus} is the action noun of \emph{spīrō} \textquotedblleft{}to breathe\textquotedblright{} (Lewis \& Short; EtymDB). English \emph{spirit} comes from it.
\end{cousinweb}

\begin{grammarlens}[title={What you've built}]
One more word for telling, explaining and describing at B1.
\end{grammarlens}

\subsection*{Your turn}
\begin{itemize}
\raggedright
  \item \emph{Say it:} \emph{spīritus, spīritūs}
  \item \emph{Say it:} \emph{spīritus, spīritūs}, once more
  \item \emph{Say it:} a sentence of your own with \emph{spīritus, spīritūs}
  \item \emph{Recall:} say the Latin for to smell, then the Latin for to perceive, then say \emph{spīritus, spīritūs} again
\end{itemize}

\begin{tcolorbox}[breakable,colback=teal!4,colframe=teal!35!black,title={Before you move on}]
What does spīritus, spīritūs mean? (\textquotedblleft{}Breath; spirit\textquotedblright{}.) Say it once more.
\end{tcolorbox}
\section[anhēlō, anhēlāre]{anhēlō, anhēlāre — to pant, gasp}

\label{lesson:LA-C520-anhelo}



\begin{quote}
Before the new one: say the Latin for to perceive, then the Latin for breath.
\end{quote}

\subsection*{You'll want to know: anhēlō, anhēlāre}
\textbf{anhēlō, anhēlāre} — \textquotedblleft{}to pant, gasp\textquotedblright{}.

\begin{quote}
Canis post cursum anhēlat.
\end{quote}

\begin{quote}
\emph{The dog pants after the run.}
\end{quote}

\emph{anhēlō, anhēlāre, anhēlāvī, anhēlātum}.

\begin{cousinweb}
\textbf{anhēlō} is built, as Lewis \& Short explains it, on a prefix \emph{an-} and \emph{hālō} \textquotedblleft{}to breathe out\textquotedblright{}, but the details are not settled.
\end{cousinweb}

\begin{grammarlens}[title={What you've built}]
One more word for telling, explaining and describing at B1.
\end{grammarlens}

\subsection*{Your turn}
\begin{itemize}
\raggedright
  \item \emph{Say it:} \emph{anhēlō, anhēlāre}
  \item \emph{Say it:} \emph{anhēlō, anhēlāre}, once more
  \item \emph{Say it:} a sentence of your own with \emph{anhēlō, anhēlāre}
  \item \emph{Recall:} say the Latin for to perceive, then the Latin for breath, then say \emph{anhēlō, anhēlāre} again
\end{itemize}

\begin{tcolorbox}[breakable,colback=teal!4,colframe=teal!35!black,title={Before you move on}]
What does anhēlō, anhēlāre mean? (\textquotedblleft{}To pant, gasp\textquotedblright{}.) Say it once more.
\end{tcolorbox}
\section[singultus, singultūs]{singultus, singultūs — a sob; a hiccup}

\label{lesson:LA-C520-singultus}



\begin{quote}
Before the new one: say the Latin for breath, then the Latin for to pant.
\end{quote}

\subsection*{You'll want to know: singultus, singultūs}
\textbf{singultus, singultūs} — \textquotedblleft{}a sob; a hiccup\textquotedblright{}.

\begin{quote}
Puella inter singultūs loquitur.
\end{quote}

\begin{quote}
\emph{The girl speaks between sobs.}
\end{quote}

Masculine, fourth declension. Breath broken by weeping, or a hiccup.

\begin{cousinweb}
\textbf{singultus} is usually traced to \emph{singulī} \textquotedblleft{}one at a time\textquotedblright{} (Lewis \& Short): speech broken into single sounds.
\end{cousinweb}

\begin{grammarlens}[title={What you've built}]
One more word for telling, explaining and describing at B1.
\end{grammarlens}

\subsection*{Your turn}
\begin{itemize}
\raggedright
  \item \emph{Say it:} \emph{singultus, singultūs}
  \item \emph{Say it:} \emph{singultus, singultūs}, once more
  \item \emph{Say it:} a sentence of your own with \emph{singultus, singultūs}
  \item \emph{Recall:} say the Latin for breath, then the Latin for to pant, then say \emph{singultus, singultūs} again
\end{itemize}

\begin{tcolorbox}[breakable,colback=teal!4,colframe=teal!35!black,title={Before you move on}]
What does singultus, singultūs mean? (\textquotedblleft{}A sob; a hiccup\textquotedblright{}.) Say it once more.
\end{tcolorbox}
\section[oscitō, oscitāre]{oscitō, oscitāre — to yawn, gape}

\label{lesson:LA-C520-oscito}



\begin{quote}
Before the new one: say the Latin for to pant, then the Latin for a sob.
\end{quote}

\subsection*{You'll want to know: oscitō, oscitāre}
\textbf{oscitō, oscitāre} — \textquotedblleft{}to yawn, gape\textquotedblright{}.

\begin{quote}
Discipulus fessus in scholā oscitat.
\end{quote}

\begin{quote}
\emph{The tired pupil yawns at school.}
\end{quote}

\emph{oscitō, oscitāre}: to open the mouth wide; also to be drowsy and careless.

\begin{cousinweb}
\textbf{oscitō} is built on \emph{ōs} \textquotedblleft{}mouth\textquotedblright{} and \emph{citō} \textquotedblleft{}to set in motion\textquotedblright{} (EtymDB): to set the mouth moving. English \emph{oscitation} \textquotedblleft{}yawning\textquotedblright{} is a rare learned word from it.
\end{cousinweb}

\begin{grammarlens}[title={What you've built}]
One more word for telling, explaining and describing at B1.
\end{grammarlens}

\subsection*{Your turn}
\begin{itemize}
\raggedright
  \item \emph{Say it:} \emph{oscitō, oscitāre}
  \item \emph{Say it:} \emph{oscitō, oscitāre}, once more
  \item \emph{Say it:} a sentence of your own with \emph{oscitō, oscitāre}
  \item \emph{Recall:} say the Latin for to pant, then the Latin for a sob, then say \emph{oscitō, oscitāre} again
\end{itemize}

\begin{tcolorbox}[breakable,colback=teal!4,colframe=teal!35!black,title={Before you move on}]
What does oscitō, oscitāre mean? (\textquotedblleft{}To yawn, gape\textquotedblright{}.) Say it once more.
\end{tcolorbox}
\section[sternuō, sternuere]{sternuō, sternuere — to sneeze}

\label{lesson:LA-C520-sternuo}



\begin{quote}
Before the new one: say the Latin for a sob, then the Latin for to yawn.
\end{quote}

\subsection*{You'll want to know: sternuō, sternuere}
\textbf{sternuō, sternuere} — \textquotedblleft{}to sneeze\textquotedblright{}.

\begin{quote}
Puer propter pulverem sternuit.
\end{quote}

\begin{quote}
\emph{The boy sneezes because of the dust.}
\end{quote}

\emph{sternuō, sternuere, sternuī}. To Roman eyes a sneeze could be an omen.

\begin{cousinweb}
\textbf{sternuō} is a native Latin verb; Lewis \& Short compares Greek \emph{ptarnusthai} \textquotedblleft{}to sneeze\textquotedblright{}, but EtymDB gives no parent, so the origin is not settled.
\end{cousinweb}

\begin{grammarlens}[title={What you've built}]
One more word for telling, explaining and describing at B1.
\end{grammarlens}

\subsection*{Your turn}
\begin{itemize}
\raggedright
  \item \emph{Say it:} \emph{sternuō, sternuere}
  \item \emph{Say it:} \emph{sternuō, sternuere}, once more
  \item \emph{Say it:} a sentence of your own with \emph{sternuō, sternuere}
  \item \emph{Recall:} say the Latin for a sob, then the Latin for to yawn, then say \emph{sternuō, sternuere} again
\end{itemize}

\begin{tcolorbox}[breakable,colback=teal!4,colframe=teal!35!black,title={Before you move on}]
What does sternuō, sternuere mean? (\textquotedblleft{}To sneeze\textquotedblright{}.) Say it once more.
\end{tcolorbox}
\section[(dialogue)]{Words from chapters 513–520 — a review}

\label{lesson:LA-R520-review-513-520}



\begin{quote}
Say the Latin for compassionate and pity.
\end{quote}

\subsection*{Your turn}
Say these aloud:

\begin{itemize}
\raggedright
  \item compassionate, pity, to pity, it shames, shyness
  \item to despise, to reject, to look down on, indignation, to displease
  \item face, appearance, hair, eyebrow, eyelid
  \item throat, jaws, neck, a limb, joints
  \item a vein, a sinew, a hide, sweat, a beat
  \item side, hip, heel, sole of the foot, navel
  \item sight, hearing, touch, to smell, to perceive
  \item breath, to pant, a sob, to yawn, to sneeze
\end{itemize}

\begin{tcolorbox}[breakable,colback=teal!4,colframe=teal!35!black,title={Before you move on}]
Compassionate? (\textbf{misericors, misericordis}.) Pity? (\textbf{misericordia, misericordiae}.) To pity? (\textbf{misereor, miserērī}.) It shames? (\textbf{pudet, pudēre}.) Shyness? (\textbf{verēcundia, verēcundiae}.) To despise? (\textbf{contemnō, contemnere}.) To reject? (\textbf{spernō, spernere}.) To look down on? (\textbf{dēspiciō, dēspicere}.) Indignation? (\textbf{indignātiō, indignātiōnis}.) To displease? (\textbf{displiceō, displicēre}.) Face? (\textbf{vultus, vultūs}.) Appearance? (\textbf{faciēs, faciēī}.) Hair? (\textbf{crīnis, crīnis}.) Eyebrow? (\textbf{supercilium, superciliī}.) Eyelid? (\textbf{palpebra, palpebrae}.) Throat? (\textbf{guttur, gutturis}.) Jaws? (\textbf{faucēs, faucium}.) Neck? (\textbf{cervīx, cervīcis}.) A limb? (\textbf{membrum, membrī}.) Joints? (\textbf{artus, artuum}.) A vein? (\textbf{vēna, vēnae}.) A sinew? (\textbf{nervus, nervī}.) A hide? (\textbf{pellis, pellis}.) Sweat? (\textbf{sūdor, sūdōris}.) A beat? (\textbf{pulsus, pulsūs}.) Side? (\textbf{latus, lateris}.) Hip? (\textbf{coxa, coxae}.) Heel? (\textbf{calx, calcis}.) Sole of the foot? (\textbf{planta, plantae}.) Navel? (\textbf{umbilīcus, umbilīcī}.) Sight? (\textbf{vīsus, vīsūs}.) Hearing? (\textbf{audītus, audītūs}.) Touch? (\textbf{tāctus, tāctūs}.) To smell? (\textbf{olfaciō, olfacere}.) To perceive? (\textbf{percipiō, percipere}.) Breath? (\textbf{spīritus, spīritūs}.) To pant? (\textbf{anhēlō, anhēlāre}.) A sob? (\textbf{singultus, singultūs}.) To yawn? (\textbf{oscitō, oscitāre}.) To sneeze? (\textbf{sternuō, sternuere}.)
\end{tcolorbox}
